\uuid{ZAyM}
\titre{Primitive}
\auteur{bodin ; croizat ; sacre}
\organisation{exo7}
\licence{CC BY-NC-SA 4.0}
\datecreate{2022-03-08}
\niveau{L1}
\module{Analyse}
\chapitre{Calcul d'intégrales}
\sousChapitre{Primitives diverses}
\difficulte{1}
\theme{Primitive}
\exo7qcmid{hilisit.equations-differentielles.003}

\contenu{
\begin{question}
\qclassification{127.11}
Parmi les phrases suivantes, quelles sont les affirmations correctes ?
\begin{answers}
  \good{L'opération du calcul de primitives est le contraire de l'opération du calcul de dérivées.}
  \good{L'opération du calcul de dérivées est le contraire de l'opération du calcul de primitives.}
  \good{Deux primitives d'une même fonction sur un intervalle sont égales à une constante près.}
  \good{Si on connaît une primitive d'une fonction, alors on les connaît toutes.}
\end{answers}
\begin{explanations}
Tout est vrai ! Les calculs de dérivées et de primitives sont bien réciproques l'un de l'autre, et dès que l'on connaît une primitive \(F\) d'une fonction \(f\) sur un intervalle, alors toutes les primitives de \(f\) sur cet intervalle seront de la forme $F(x) + C$ (où $C$ est une constante).
\end{explanations}
\end{question}
}
